%% ============================================================
%% Appendix D: Periodic Solutions and Floquet Theory
%% ============================================================

\chapter{Periodic Solutions and Floquet Theory}
\label{app:floquet}

For models with periodic forcing or limit-cycle behavior, SA of the
periodic orbit requires tools beyond those developed in Chapter~4.
This appendix introduces Floquet theory and the sensitivity of
periodic solutions, based on the 2009 monograph~\cite{arriola2009sensitivity}.

%%------------------------------------------------------------
\section{Floquet Theory}
\label{appsec:floquet}
%%------------------------------------------------------------

Consider the linear, $T$-periodic ODE $\dot{\vec{u}} = A(t)\vec{u}$
where $A(t+T) = A(t)$. The \textbf{fundamental matrix solution}
$\Phi(t)$ satisfies $\dot{\Phi} = A(t)\Phi$, $\Phi(0) = I$.

The \textbf{monodromy matrix}\index{monodromy matrix} is $M = \Phi(T)$: the fundamental
matrix evaluated at the end of one period. Its eigenvalues
$\mu_1, \ldots, \mu_n$ are the \textbf{Floquet multipliers}.\index{Floquet multipliers}

\textbf{Floquet's Theorem.}\index{Floquet theory} The fundamental matrix has the form
$\Phi(t) = P(t)e^{Bt}$ where $P(t+T) = P(t)$ is periodic and
$B$ is a constant matrix satisfying $e^{BT} = M$.

The Floquet multipliers determine stability:
if all $|\mu_k| < 1$ (except the trivial multiplier $\mu = 1$
corresponding to time translation), the periodic orbit is stable.
If any $|\mu_k| > 1$, the orbit is unstable.

%%------------------------------------------------------------
\section{Sensitivity of the Period}
\label{appsec:periodsens}
%%------------------------------------------------------------

For a nonlinear system $\dot{\vec{u}} = \vec{F}(\vec{u}; \vec{p})$
with a periodic orbit $\vec{u}^*(t; \vec{p})$ of period $T(\vec{p})$,
the period itself is a \QOI. Its sensitivity with respect to parameter
$p_j$ is:
\begin{equation}
  \dd{T}{p_j} = -\frac{\vec{s}^T\pd{\vec{F}}{p_j}(\vec{u}^*(T), \vec{p})}
                       {\vec{s}^T\vec{F}(\vec{u}^*(T), \vec{p})},
  \label{eq:periodsens}
\end{equation}
where $\vec{s}$ is the left eigenvector of $M$ corresponding to
the unit eigenvalue ($M^T\vec{s} = \vec{s}$, $\|\vec{s}\| = 1$).

This formula is derived by differentiating the period-map condition
$\vec{u}^*(0) = \vec{u}^*(T)$ with respect to $p_j$, accounting for
the fact that the period $T$ itself changes.

%%------------------------------------------------------------
\section{The van der Pol Example}
\label{appsec:vanderpol}
%%------------------------------------------------------------

The van der Pol oscillator $\ddot{u} - \epsilon(1-u^2)\dot{u} + u = 0$
(written as a system: $\dot{u}_1 = u_2$, $\dot{u}_2 = \epsilon(1-u_1^2)u_2 - u_1$)
has a stable limit cycle for $\epsilon > 0$.

For small $\epsilon$, the period is approximately $T \approx 2\pi(1 + \epsilon^2/16)$.
The period sensitivity is $\SI{\epsilon}{T} \approx \epsilon^2/(8\pi)$,
which is small for small $\epsilon$ but grows with increasing nonlinearity.

For larger $\epsilon$, the limit cycle becomes increasingly asymmetric
and the period sensitivity must be computed numerically using the
monodromy matrix approach described above.

%%------------------------------------------------------------
\section{Sensitivity of Floquet Multipliers}
\label{appsec:floquetsens}
%%------------------------------------------------------------

The sensitivity of the Floquet multiplier $\mu_k$ with respect to
parameter $p_j$ uses the eigenvalue sensitivity formula from
Appendix~\ref{appsec:eigensens} applied to the monodromy matrix:
\begin{equation}
  \dd{\mu_k}{p_j} = \frac{\vec{y}_k^T\pd{M}{p_j}\vec{x}_k}{\vec{y}_k^T\vec{x}_k},
  \label{eq:floquetsens}
\end{equation}
where $\vec{x}_k$ and $\vec{y}_k$ are the right and left eigenvectors
of $M$ for eigenvalue $\mu_k$, and $\partial M/\partial p_j$ is computed
by integrating the variational equation for $\partial\Phi/\partial p_j$
over one period.

When $|\mu_k| \to 1$ for $k \neq 1$, the periodic orbit approaches
a bifurcation (Hopf, period-doubling, or Neimark-Sacker), and the
sensitivity of $\mu_k$ diverges, exactly analogous to the
divergence of SI near a transcritical bifurcation discussed in Chapter~8.

%%------------------------------------------------------------
\section{Exercises}
\label{appsec:D_exercises}
%%------------------------------------------------------------

\begin{selfcheck}
The van der Pol oscillator with $\epsilon = 0.5$ has a stable limit cycle.
(a) Using the approximation $T \approx 2\pi(1 + \epsilon^2/16)$,
compute the period and the normalized sensitivity $\SI{\epsilon}{T}$.
(b) Is the period more or less sensitive to $\epsilon$ than a linear oscillator's
period is to its natural frequency? Explain.
(c) As $\epsilon$ increases from $0.5$ to $2.0$, does the period sensitivity
increase or decrease? What does this imply about using local SA at
large $\epsilon$?
\end{selfcheck}
\begin{scAnswer}
(a) $T \approx 2\pi(1 + 0.25/16) = 2\pi \times 1.01563 \approx 6.377$.
$dT/d\epsilon = 2\pi(\epsilon/8) = 2\pi(0.0625) \approx 0.393$.
$\SI{\epsilon}{T} = (\epsilon/T)(dT/d\epsilon) = (0.5/6.377)(0.393) \approx 0.031$.
(b) A linear oscillator's period is independent of amplitude (and $\epsilon$ corresponds
to the nonlinear damping coefficient, not the natural frequency). For $\epsilon=0.5$,
the period sensitivity $\approx 0.031$ is small: the period is nearly insensitive
to $\epsilon$ in the weakly nonlinear regime.
(c) As $\epsilon$ increases, $d\lambda_k/d\varepsilon$ grows because the limit cycle
becomes more asymmetric and the monodromy matrix changes more rapidly with $\epsilon$.
The local SI underestimates the true sensitivity at large $\epsilon$, which is precisely
the scenario where global SA (Chapter~9) is needed.
\end{scAnswer}

\begin{selfcheck}
A compartmental epidemic model exhibits a periodic outbreak pattern with
period $T^* = 365$~days (annual epidemic) under seasonal forcing.
The Floquet multiplier associated with the periodic orbit is $\mu = 0.85$.
(a) Is this periodic orbit stable or unstable? How do you know?
(b) If a parameter perturbation changes $\mu$ to $1.02$, what happens
to the epidemic dynamics?
(c) The sensitivity of $\mu$ to the transmission rate $\beta$ is
$d\mu/d\beta = 3.2$ (computed via~\eqref{eq:floquetsens}).
If $\beta$ increases by $5\%$, estimate the new value of $\mu$.
Is the periodic outbreak still stable?
\end{selfcheck}
\begin{scAnswer}
(a) Stable: $|\mu| = 0.85 < 1$, so perturbations from the periodic orbit
decay over successive periods. The orbit is an attractor.
(b) If $\mu = 1.02 > 1$, the periodic orbit becomes unstable: nearby
trajectories drift away from the annual cycle. The epidemic may lose its
annual periodicity or transition to a different attractor (period-doubling,
chaos, or endemic equilibrium).
(c) $\Delta\mu \approx (d\mu/d\beta)\,\Delta\beta = 3.2\times(0.05\times\beta)$.
If $\beta^* = 0.3$ (typical), $\Delta\mu \approx 3.2\times0.015 = 0.048$.
New $\mu \approx 0.85 + 0.048 = 0.898 < 1$: still stable.
The first-order approximation suggests stability is preserved for this perturbation size.
\end{scAnswer}
